% All-order free detector used by trace calibration and the finite scanner.
\section{A positive free detector at every jet order}
\label{sec:ex-lattice-all-grade-free-detectors}

This section proves the all-grade reference-field separation needed in
Section~\ref{sec:ex-lattice-firesky-gallery} and for the finite point
scanner. The point is that positivity must hold for a mixture of derivative
orders, not just separately at each order. We establish it by an explicit
triangular change of probes. All vector spaces in the one-leg calculation
are real; ordinary complex conjugation is added only when passing to the
Wick-polynomial pairing.

\subsection{The Kelvin response and a triangular family of probes}

Let $\mathscr S_j$ be the finite-dimensional real space of homogeneous
degree-$j$ Maxwell field strengths in four-dimensional Euclidean space,
with the fixed gauge-color representation. Thus $P\in\mathscr S_j$ is an
antisymmetric two-tensor-valued homogeneous polynomial satisfying the
Maxwell differential and algebraic identities. Let $\langle\cdot,\cdot\rangle_F$
denote the positive Frobenius inner product, including the positive color
inner product. For a nonzero point $x$, write
\[
 r_x=|x|^{-2},\qquad I(x)=r_xx,\qquad
 \vartheta_x=I-2r_x\,x\otimes x.
\]
The matrix $\vartheta_x$ is the orthogonal reflection that reverses the
line through $x$ and fixes its orthogonal complement. Its induced action on
two-tensors is denoted $\vartheta_x^{(2)}$.

For a symmetric coefficient tensor $A$ of order $j$ with two antisymmetric
field indices, write componentwise
\[
 (\mathcal D_A f)_{im}(x)=\sum_{(a_1,\ldots,a_j)\in\{0,1,2,3\}^j}
   \frac{1}{j!}D^jf(x)[e_{a_1},\ldots,e_{a_j}]\,
       A_{im}(e_{a_1},\ldots,e_{a_j}).
\]
The sum is over ordered derivative tuples; later the two field indices are
summed in both orders in the Frobenius pairing. This is the normalized
Taylor coefficient convention. The colour pairing is the fixed positive
real inner product. Write $c_j=(-2)^j/(4\pi^2)$, the coefficient of the
four-dimensional Green function.

We first establish the differentiated Green identity used below:
\begin{equation}\label{eq:ex-lattice-shell-green}
 (\mathcal D_A G)_{im}(x)=c_j|x|^{-2(j+1)}A_{im}(x),\qquad x\ne0,
 \quad A\in\mathscr S_j,
 \qquad G(x)=\frac{1}{4\pi^2|x|^2}.
\end{equation}
For $j=0$ this is evaluation of the Green function. Suppose it holds at
order $j$ and let $A$ have order $j+1$. For each coordinate direction
$e_a$, freeze the first tensor slot to get $A_a\in\mathscr S_j$. The
normalized Taylor coefficient satisfies
\[
 (\mathcal D_A f)_{im}=\frac1{j+1}\sum_a\partial_a(\mathcal D_{A_a}f)_{im}.
\]
This follows by expanding the ordered derivative tuples: each of the
$j+1$ possible first slots gives the same symmetric coefficient. Applying
the induction hypothesis and differentiating gives
\[
 (\mathcal D_A G)_{im}=\frac{c_j}{j+1}\left(
  -2(j+1)|x|^{-2(j+2)}\sum_a x_a(A_a)(x)
  +|x|^{-2(j+1)}\sum_a\partial_a(A_a)(x)\right).
\]
For completeness, the trace identity here is harmonicity in the derivative
slots, not the Maxwell divergence equation alone. Write $F_{im}(x)$ for
the homogeneous field components represented by $A$. The Maxwell equations
give $\partial_aF_{ai}=0$ and
$\partial_aF_{im}+\partial_iF_{ma}+\partial_mF_{ai}=0$. Apply $\partial_a$
and sum over $a$ in the second identity. The two terms involving the
divergence vanish, so $\Delta F_{im}=0$. By the normalized coefficient
convention, $(A_a)_{im}=(j+1)^{-1}\partial_aF_{im}$; hence
\[
 \sum_a\partial_a(A_a)_{im}
   =(j+1)^{-1}\Delta F_{im}=0.
\]
Euler's homogeneous identity, or direct contraction of the first slot,
gives
\[
 \sum_a x_a(A_a)_{im}(x)=F_{im}(x)=A_{im}(x).
\]
Since $c_{j+1}=-2c_j$, this proves
\eqref{eq:ex-lattice-shell-green}. It is an off-origin identity; no
distributional contact term is involved.

The Euclidean Green kernel and its two curls give a response field
associated with $P$. The Kelvin calculation needed here is the following
identity, with the normalization fixed by the Green function:
\begin{equation}\label{eq:ex-lattice-kelvin-response}
 \mathcal R(P)(y)
 =\operatorname{strength}(j)|y|^{-4}
   \vartheta_y^{(2)}P(I(y)),\qquad y\ne0,
 \quad\operatorname{strength}(j)
 =4(j+1)(j+2)(-2)^j\operatorname{normal}.
\end{equation}
Here $\operatorname{normal}=1/(4\pi^2)>0$ is the scalar in
\eqref{eq:ex-lattice-shell-green}. In particular,
\begin{equation}\label{eq:ex-lattice-strength-sign}
 (-1)^j\operatorname{strength}(j)>0.
\end{equation}
The sign follows directly from the displayed formula: the factors
$4(j+1)(j+2)$ and $\operatorname{normal}$ are positive and
$(-1)^j(-2)^j=2^j$.

We recall the differential reduction giving this normalization. Put
$R_i(x)=\sum_a x_aP_{ai}(x)$ and
$c_j=(-2)^j\operatorname{normal}$. The differentiated Green potential is
$c_j|x|^{-2(j+1)}P_{ai}(x)$. The divergence-free equation removes the
derivatives falling on $P$ in its divergence, so the resulting vector is
\[
 D_i(x)=-2(j+1)c_j|x|^{-2(j+2)}R_i(x).
\]
The curl of the field response is
$2(\partial_mD_i-\partial_iD_m)$. Closedness, the Bianchi identity, and
Euler homogeneity give
\[
 \partial_iR_m(x)-\partial_mR_i(x)=(j+2)P_{im}(x).
\]
Also $\partial_i|x|^{-2(j+2)}=-2(j+2)x_i|x|^{-2(j+3)}$. Substitution
yields
\[
 4(j+1)(j+2)c_j|x|^{-2(j+2)}
 \left[
   P_{im}(x)-2|x|^{-2}\bigl(x_iR_m(x)-x_mR_i(x)\bigr)
 \right].
\]
The bracket is the two-index reflection of $P(x)$. Homogeneity gives
\[
 |x|^{-2(j+2)}\vartheta_x^{(2)}P(x)
   =|x|^{-4}\vartheta_x^{(2)}P(I(x)).
\]
This proves
(\ref{eq:ex-lattice-kelvin-response}) off the origin and identifies its
coefficient; no contact term is used.

For an orthogonal reflection $s$, define its full action on a solution by
$(s\cdot Q)(v)=s^{(2)}Q(sv)$: it acts on both field indices and every
derivative slot. For $x\ne0$, let $\operatorname{Leg}_x(P,Q)$ be the
one-leg kernel: the covariance of the probe $P$ at a point $y_0$ with the
probe $Q$ at $y_0-x$, that is, the second derivatives of the scalar Green
kernel at separation $x$ (first point minus second point) contracted with
the dual probes of $P$ and $Q$, followed by the two curls.  The result does
not depend on $y_0$. The one-leg reflected pairing is
\[
 B_x^{j,k}(P,Q)=\operatorname{Leg}_x(P,\vartheta_x\cdot Q).
\]
Orthogonal covariance and exchange symmetry of the Green kernel imply
\[
 B_x^{j,k}(P,Q)=B_x^{k,j}(Q,P).
\]
The use of the reflection in the second slot is essential: this is the
one-leg reflected form, not the unreflected covariance.

To triangularize all grades at once, for $0\leq k\leq j$ let
\[
 T_{j,k}^{x}P
 =\frac{\operatorname{strength}(j)}{\operatorname{strength}(k)}
   \operatorname{Taylor}_{k}\bigl(P(\,\cdot-I(x))\bigr),
 \qquad P\in\mathscr S_j.
\]
The Taylor coefficient is taken at the origin, and the translated
coefficient is again a homogeneous degree-$k$ Maxwell solution. Set
\[
 \Sigma_xP=\sum_{k=0}^{j} T_{j,k}^{x}P\in\mathscr S_{<j+1},
 \qquad P\in\mathscr S_j.
\]
The sum is finite. By the exact Taylor formula for a degree-$j$ polynomial,
for every $y$,
\begin{equation}\label{eq:ex-lattice-shift-taylor}
 \sum_{k=0}^{j}\operatorname{strength}(k)
   (T_{j,k}^{x}P)(y)
 =\operatorname{strength}(j)P(y-I(x)).
\end{equation}
Apply (\ref{eq:ex-lattice-kelvin-response}) term by term and use
(\ref{eq:ex-lattice-shift-taylor}). The resulting response is
$\mathcal R_x(\Sigma_xP):=\sum_{k=0}^j\mathcal R(T_{j,k}^{x}P)$.
For $y$ near $x$, inversion gives
\[
I(x+t v)-I(x)=t\,V_{x,v}(t),\qquad
V_{x,v}(t)=|x+t v|^{-2}
  \bigl(\vartheta_xv-t|v|^2I(x)\bigr),
\]
where $V_{x,v}$ is smooth near $t=0$. Since $P$ is homogeneous of degree
$j$, the shifted response along this line is $t^j$ times a smooth function.
Consequently its Taylor coefficients through order $j-1$ vanish, while
its order-$j$ coefficient is
\[
 \operatorname{strength}(j)|x|^{-2(j+2)}
  (\vartheta_x\cdot P)(v),
\]
where the reflected solution is defined by
$(\vartheta_x\cdot P)(v)=\vartheta_x^{(2)}P(\vartheta_xv)$.
To connect the response computation to the one-leg pairing, the
two-curl Green identity gives
\[
 B_x(\Sigma_xP,Q_k)=(-1)^k
 \left\langle \vartheta_x\cdot Q_k,
   \operatorname{Taylor}_{k}\bigl(\mathcal R_x(\Sigma_xP)\bigr)(x)\right\rangle_F.
\]
Here $Q_k\in\mathscr S_k$, and the normalized Taylor tensor includes all
ordered derivative tuples. The sign $(-1)^k$ is the sign from
differentiating the second insertion, which sits at $y_0-x$: the pulse of a
grade-$k$ probe is odd or even under $y\mapsto-y$ according to the parity of
$k$.  The triangular transform below uses this orientation; with the
opposite one the sign would attach to the components of the first argument. The vanishing lower Taylor
coefficients imply $B_x(\Sigma_xP,Q_k)=0$ when $k<j$; at $k=j$ the
order-$j$ coefficient gives
$B_x(\Sigma_xP,Q_j)=\lambda_j(x)\langle P,Q_j\rangle_F$, because
$(\vartheta_x\cdot Q_j,\vartheta_x\cdot P)$ have the same Frobenius
product as $(Q_j,P)$. When the second input is also transformed, its
components have grades at most its original grade, and its top component is
unchanged because $T_{k,k}^x$ is the identity. Thus these same two
identities give the triangular identities for the full one-leg kernel.
Extend $B_x$ by bilinearity to the finite tower. If
$\mathscr S_{<m}=\bigoplus_{k<m}\mathscr S_k$, then the formulas
\begin{align}
 B_x(\Sigma_xP,\Sigma_xQ)&=0
     &&(P\in\mathscr S_j,\ Q\in\mathscr S_k,\ k<j),
       \label{eq:ex-lattice-shift-cross}\\
 B_x(\Sigma_xP,\Sigma_xQ)&=\lambda_j(x)
     \langle P,Q\rangle_F
     &&(P,Q\in\mathscr S_j),
       \label{eq:ex-lattice-shift-diagonal}
\end{align}
where
\begin{equation}\label{eq:ex-lattice-positive-pivot}
 \lambda_j(x)=(-1)^j\operatorname{strength}(j)|x|^{-2(j+2)}>0.
\end{equation}
The factor $(-1)^k$ in the one-leg kernel comes from differentiating the
second insertion; at equal grade it combines with the Taylor coefficient
to give precisely (\ref{eq:ex-lattice-positive-pivot}).

Now fix a finite cutoff $n$ and form the tower
\[
 \mathscr S_{<n}=\bigoplus_{j=0}^{n-1}\mathscr S_j.
\]
On this tower the reflected Gram form is
\[
 B_x(P,Q)=\sum_{j,k<n}B_x^{j,k}(P_j,Q_k).
\]
Apply the triangular transform $\Sigma_x$ separately in each input grade.
The cross-grade identities (\ref{eq:ex-lattice-shift-cross}) and their
transposes remove every off-diagonal block, while
(\ref{eq:ex-lattice-shift-diagonal}) gives
\begin{equation}\label{eq:ex-lattice-tower-gram}
 B_x(\Sigma_xP,\Sigma_xQ)
 =\sum_{j<n}\lambda_j(x)\langle P_j,Q_j\rangle_F.
\end{equation}
Every coefficient on the right is positive. Since each Frobenius inner
product is positive definite, the right side is positive for every nonzero
tower. Moreover, $T_{j,j}^{x}$ is the identity: the top-degree Taylor
coefficient of $P(\,\cdot-I(x))$ is $P$. Therefore the matrix of $\Sigma_x$
on the ordered grades is lower triangular with identity diagonal. Its
strictly grade-lowering part $N_x=\Sigma_x-I$ satisfies $N_x^n=0$, so
\[
 \Sigma_x^{-1}=I-N_x+N_x^2-\cdots+(-1)^{n-1}N_x^{n-1}.
\]
Thus the diagonalization is an invertible change of probes, and the full
mixed-grade Gram form is positive definite. This is a
simultaneous statement about the entire tower; positivity of the separate
diagonal blocks alone would not imply it.

\subsection{Wick polynomials and strict reflected positivity}

Let $\mathscr A$ be the complex polynomial algebra generated by the formal
field-strength jet coordinates at all derivative orders. Evaluation on
homogeneous Maxwell solutions restricts each polynomial to the on-shell
jet space, imposing both the differential Maxwell equations and the
algebraic Bianchi identities. The filtered shell sector $\mathscr Q_{\le d}$
is the image of this restriction map through grade $d$, hence is canonically
the quotient by its kernel. Full reflection preserves each grade and this
filtered sector. Each linear coordinate restricted to $\mathscr S_j$ is a
linear functional on a finite-dimensional inner-product space, so the
finite-dimensional Riesz representation theorem writes it as pairing with
a vector of $\mathscr S_j$. Consequently every polynomial in finitely many
jet coordinates is a polynomial in the coordinates of one sufficiently
large finite tower. A homogeneous solution $P\in\mathscr S_j$
defines its linear Wick symbol $\mathsf d(P)$. Products of these symbols
generate the polynomial Wick algebra.

For two distinct points $x,y$, put $z=x-y$. Let
$\mathcal W_{x,y}(A,B)$ denote the free two-point moment of the polynomial
insertions $A$ and $B$. The reflected pairing is linear in its first
argument and conjugate-linear in its second:
\[
\langle A,B\rangle_{x,y}^{\mathrm{ref}}
   =\mathcal W_{x,y}\bigl(A,\vartheta_z\cdot\overline{B}\bigr),
\]
where $\vartheta_z$ acts on every tensor index of the second field
insertion. Complex conjugation of polynomial coefficients is a separate
operation in this formula. The free moment is the normal-ordered cross-site
Gaussian matching moment: only pairings joining the two distinct insertion
sites occur. Removing one linear field factor from the first insertion
therefore gives the sum of its contractions with the factors in the second
insertion. This follows by separating the unique partner of the removed
factor in the finite matching formula for a free moment. With two insertion
points there is only one possible partner site. In orthogonal coordinates
the identity reads
\begin{equation}\label{eq:ex-lattice-script-recurrence}
 \langle X_i p,q\rangle_{x,y}^{\mathrm{ref}}
   =d_i\langle p,\partial_iq\rangle_{x,y}^{\mathrm{ref}},
 \qquad
 \langle1,X_iq\rangle_{x,y}^{\mathrm{ref}}=0,
 \qquad \langle1,1\rangle_{x,y}^{\mathrm{ref}}=1.
\end{equation}
The first identity is the differentiated removal rule, and the second is
the same rule with no factor available to contract on the first side.

Choose a real orthogonal basis $e_1,\ldots,e_M$ for the tower Gram form
\eqref{eq:ex-lattice-tower-gram} at the separation $z=x-y$, with squared
norms $d_i>0$. Such a basis
exists by finite-dimensional Gram--Schmidt: at each step, the orthogonal
complement of the span already chosen is positive definite, so a nonzero
vector in it has strictly positive squared norm; normalize or retain its
positive norm. In the polynomial coordinates $X_1,\ldots,X_M$ associated
with $e_i$, the reflected pairing satisfies
\[
 \langle X^\alpha,X^\beta\rangle_{x,y}^{\mathrm{ref}}
 =\begin{cases}
   \displaystyle\prod_{i=1}^M \alpha_i!d_i^{\alpha_i},&\alpha=\beta,\\
   0,&\alpha\ne\beta.
  \end{cases}
\]
For completeness, this formula follows by induction on $|\alpha|$. When
$\alpha=0$, the last two identities in
(\ref{eq:ex-lattice-script-recurrence}) give $1$ if $\beta=0$ and $0$
otherwise. If $\alpha_i>0$, apply the first identity with
$p=X^{\alpha-e_i}$ and $q=X^\beta$:
\[
 \langle X^\alpha,X^\beta\rangle
   =\beta_i d_i\langle X^{\alpha-e_i},X^{\beta-e_i}\rangle,
\]
where the right side is zero when $\beta_i=0$. Induction shows that the
pairing vanishes unless $\alpha=\beta$. When they agree, removing each
occurrence of coordinate $i$ contributes successively
$d_i\alpha_i,d_i(\alpha_i-1),\ldots,d_i$, so its total factor is
$\alpha_i!d_i^{\alpha_i}$. Multiplication over coordinates gives the
formula. For arbitrary polynomials
$p=\sum_\alpha p_\alpha X^\alpha$ and
$q=\sum_\alpha q_\alpha X^\alpha$, it follows that
\begin{equation}\label{eq:ex-lattice-wick-norm}
 \langle p,q\rangle_{x,y}^{\mathrm{ref}}
 =\sum_{\alpha}p_\alpha\overline{q_\alpha}
     \prod_{i=1}^M\alpha_i!d_i^{\alpha_i},
 \qquad
 \langle p,p\rangle_{x,y}^{\mathrm{ref}}
 =\sum_{\alpha}|p_\alpha|^2
     \prod_{i=1}^M\alpha_i!d_i^{\alpha_i}.
\end{equation}
Each weight is strictly positive, including the empty multi-index, whose
weight is $1$. Thus the reflected norm is strictly positive for every
nonzero polynomial. Every polynomial uses only finitely many jet orders;
choose a tower cutoff containing those orders, and the argument applies.
If two representatives define the same shell polynomial, their difference
restricts to zero on every solution and has zero pairing with every
insertion. The formula is therefore well defined on the shell quotient and
strictly positive on each nonzero class.

\begin{theorem}[Strict all-grade free detector]
\label{thm:ex-lattice-all-grade-free-detector}
For any two distinct Euclidean points, the reflected free two-point form on
the polynomial algebra of gauge-invariant field-strength jets, after
quotienting by the Maxwell equations and the Bianchi identities, is
strictly positive: if $P$ is a nonzero quotient class, then
\[
 \langle P,P\rangle_{x,y}^{\mathrm{ref}}>0.
\]
Consequently, the family of separated two-point evaluations at any fixed
pair detects every field in each finite-dimensional grade.
\end{theorem}

\begin{proof}
The Kelvin calculation gives the one-leg triangular identities and the
strictly positive tower pivots in
(\ref{eq:ex-lattice-shift-cross})--(\ref{eq:ex-lattice-positive-pivot}).
Equation (\ref{eq:ex-lattice-tower-gram}) therefore diagonalizes the
entire finite tower with positive diagonal entries. The removal-rule
calculation gives (\ref{eq:ex-lattice-wick-norm}), whose summands are
nonnegative and whose summand for any nonzero coefficient is positive.
Every polynomial lies in a finite tower, proving strict positivity for
each nonzero shell class.

For the detector statement, fix a grade $d$ and a nonzero class $v$ in its
free field space. Choose any separated pair. The reflected evaluation of
$v$ against its reflected conjugate is the strictly positive number just
proved. Hence at least one vector in any spanning basis of the grade has a
nonzero pairing with $v$; otherwise linearity would make its pairing with
every grade-$d$ vector, including the reflected conjugate of $v$, vanish.
Thus the full family of basis-column evaluations is injective.
\end{proof}

\subsection{The finite point scanner}

We spell out the map used by the construction. Fix a reference frame and a
grade $d$. Choose a basis $b_0,\ldots,b_m$ of its free field space with
$b_0=1$, the scalar unit; the remaining basis elements are images of
vectors in that same reference frame. At a separated pair $(x,y)$ define
the row scanner
\[
 \mathcal R_{x,y}(v)
   =\bigl(\operatorname{Use}_{x,y}(v,b_j)\bigr)_{j=0}^{m},
 \qquad
 \operatorname{Face}_{x,y}(w)
   =\vartheta_{x-y}\bigl(\overline w\bigr).
\]
Here the bar is ordinary coefficient conjugation, and the reflection acts
on all tensor indices. By the definitions of the free-field quotient and
the reflected polynomial pairing,
\begin{equation}\label{eq:ex-lattice-scanner-reflection}
 \operatorname{Use}_{x,y}
   (v,\operatorname{Face}_{x,y}(w))
 =\langle v,w\rangle_{x,y}^{\mathrm{ref}}.
\end{equation}
If $\mathcal R_{x,y}(v)=0$, then the bilinear form
$w\mapsto\operatorname{Use}_{x,y}(v,w)$ vanishes on the chosen basis and
therefore on the whole grade. In particular, take
$w=\operatorname{Face}_{x,y}(v)$. Equation
(\ref{eq:ex-lattice-scanner-reflection}) and strict positivity imply
$v=0$. So the scanner is injective. Notice that it is the conjugate and
the reflection together that place the second entry in the reflected
positive form; reflection alone or conjugation alone is not the required
operation.

The scalar unit and source-frame basis make this an admissible finite
scanner for the later construction. The unit is retained as the first
column, and the other columns use the same frame from which the free table
was built. This proves the all-grade reference detector only. It makes no
positivity claim about the interacting measure, and the eventual
interacting comparison and replication bounds remain separate inputs.
